\chapter{Embedding Numerics: Stratified Solvers for $\iota : \mathcal{G} \to \mathcal{M}$}
\label{ch:74}

\section{Premise}

The embedding $\iota : \mathrm{Inc}(\mathcal{G}) \hookrightarrow \mathcal{M}$ 
is not merely a mathematical abstraction.  
It is the \emph{computational bridge} through which:

\begin{itemize}
\item discrete Polyxan structure $\mathcal{G}$,
\item continuous RSVP fields $(\Phi,\mathbf{v},S)$,
\item spectral curvature $\mathcal{R}$,
\item entropy $S$,
\end{itemize}

interact during execution.

The embedding evolves by the geometric flow
\[
 \partial_t \iota
   = \Delta_{\mathrm{Inc}} \iota
     - \beta \bigl(\mathcal{R}\circ\iota - \mathcal{H}_0[\mathcal{G}]\bigr)\nabla\mathcal{R}
     - \gamma \bigl(S\circ\iota - \mathcal{H}_{\mathrm{tag}}[\mathcal{G}]\bigr)\nabla S,
\]
derived in Chapter~\ref{ch:32}.  
This chapter constructs the computational machinery for solving this flow on real hardware:

\begin{enumerate}
\item a stratified spatial discretization of $\mathrm{Inc}(\mathcal{G})$,
\item numerical evaluation of curvature and entropy mismatches,
\item stable implicit--explicit time-stepping,
\item handle-detection and topology–rewrite callbacks,
\item and harmonization with the Rewrite Scheduler (Ch.~\ref{ch:73}).
\end{enumerate}

\section{Stratified Representation of the Incidence Complex}

The incidence complex has mixed dimension:
\[
  \mathrm{Inc}(\mathcal{G}) = 
   \mathcal{V} \;\sqcup\; \mathcal{E} \;\sqcup\; \mathcal{F} \;\sqcup \cdots
\]
(vertices, hyperedges, higher cells).

For computation we treat each stratum separately.

\subsection{0–cells (vertices)}

Vertices are embedded as points:
\[
 \iota(v_i) \in \mathcal{M}.
\]

Local neighbourhood:
\[
 U_i = \{\,v_j \mid e \text{ contains } v_i,v_j\,\}.
\]

\subsection{1–cells (hyperedges)}

Each hyperedge $e$ is represented by a polyline or Bézier segment
\[
 \iota(e) = (p_0,\dots,p_k).
\]

\subsection{2–cells and higher}

Faces or higher cells are triangulated using the embedding–pullback metric:
\[
 g_{ab}^{\text{pull}} = \partial_a \iota^i\, \partial_b \iota^j \, g_{ij}.
\]

This allows Laplace–Beltrami discretization on each stratum.

\section{Discrete Laplace–Beltrami Operator on the Incidence Complex}

The operator $\Delta_{\mathrm{Inc}}$ acts differently on each stratum.

\subsection{Vertices}

Use the cotangent Laplacian:
\[
 (\Delta \iota)_i
   = \frac{1}{2A_i}
     \sum_{j\in U_i} (\cot\alpha_{ij} + \cot\beta_{ij})(\iota_j - \iota_i),
\]
with $A_i$ the mixed Voronoi area.

\subsection{Hyperedges}

Edges use a 1D Laplacian with tension term:
\[
 (\Delta_e \iota)(s)
   = \partial_s^2 \iota(s) 
     - \frac{\partial_s \|\partial_s \iota\|^2}
            {2\,\|\partial_s \iota\|^2}\,\partial_s \iota(s).
\]

\subsection{Higher cells}

On faces, apply the FEM stiffness matrix built from
$g_{ab}^{\text{pull}}$.

\section{Curvature and Entropy Mismatch Evaluation}

At each embedded point $x = \iota(\sigma)$ for cell $\sigma$:
\[
 \delta\mathcal{R}(x) := \mathcal{R}(x) - \mathcal{H}_0[\mathcal{G}],
\]
\[
 \delta S(x) := S(x) - \mathcal{H}_{\mathrm{tag}}[\mathcal{G}].
\]

Gradients are taken from the RSVP field solver (Ch.~\ref{ch:72}):
\[
 \nabla\mathcal{R}(x)
 ,\qquad 
 \nabla S(x)
\]
interpolated on the embedding.

\section{Time-Stepping the Embedding Flow}

The flow contains:
\[
 \Delta_{\mathrm{Inc}}\iota 
  \quad\text{(stiff term)},\qquad
  - \beta(\cdot)\nabla\mathcal{R}
  - \gamma(\cdot)\nabla S
 \quad\text{(non-stiff drift)}.
\]

We use an IMEX scheme:
\[
 \frac{\iota^{n+1} - \iota^n}{\Delta t}
   = \Delta_{\mathrm{Inc}} \iota^{n+1}
     + f(\iota^n),
\]
with drift
\[
 f(\iota)
   = -\beta\,\delta\mathcal{R}(\iota)\,\nabla\mathcal{R}
     -\gamma\,\delta S(\iota)\,\nabla S.
\]

Solving for $\iota^{n+1}$ requires one linear solve per stratum.

\section{Topology Change Detection}

Topology change occurs when:
\begin{enumerate}
\item two vertices satisfy 
      $\|\iota(v_i)-\iota(v_j)\| < \varepsilon_{\text{merge}}$;
\item curvature on an edge blows up  
      $\|\partial_s\iota\|^{-1} > \varepsilon_{\text{snap}}$;
\item the Jacobian of a face embedding approaches zero.
\end{enumerate}

When detected, a callback triggers an EHR rewrite with an event certificate (Ch.~\ref{ch:73}).

\section{Coupling to the Rewrite Scheduler}

Whenever rewriting occurs:
\begin{enumerate}
\item \textbf{Remove} collapsed strata from $\mathrm{Inc}(\mathcal{G})$.
\item \textbf{Insert} new strata from $R$ in the rewrite $L \leftarrow K \to R$.
\item \textbf{Reinitialize} $\iota$ on new cells via harmonic extension:
\[
 \Delta_{\mathrm{Inc}} \iota_{\text{new}} = 0
\]
with Dirichlet boundary on $K$.
\item \textbf{Continue} the IMEX solver with updated geometry.
\end{enumerate}

\section{Stability Guarantees}

\begin{theorem}[Stability of the Embedding Solver]
For $\Delta t \le C \, h^2$ (mesh size $h$),  
the IMEX scheme is unconditionally stable on each stratum  
and preserves curvature-priority descent.
\end{theorem}

\begin{theorem}[Finite-Time Regularity]
The stratified solver produces a piecewise-smooth embedding  
for all $t < T_{\mathrm{collapse}}$  
and triggers exact topology changes at $T_{\mathrm{collapse}}$.
\end{theorem}

\section{Interpretability of Embedding Updates}

Each update $\iota^{n+1}$ includes a certificate:
\[
  \mathsf{Cert}_{\iota}(n)
   := \bigl(\Delta_{\mathrm{Inc}}, f, \delta\mathcal{R}, \delta S \bigr),
\]
mirroring the classifier for discrete rewrites.
This allows the verification layer (Part~IV)
to reconstruct:
\begin{itemize}
\item the origin of each movement,
\item curvature and entropy drivers,
\item and justification for every topology change.
\end{itemize}

\section{Final Cadence}

The Embedding Solver is the computational realization of the statement:

\begin{quote}
\emph{Geometry acts on syntax.  
Syntax reacts on geometry.  
Numerics make their unity executable.}
\end{quote}

With this chapter, the Semantic Execution Engine gains its complete geometric half.
Coupled with the Rewrite Scheduler (Ch.~\ref{ch:73}),
it now evolves meaning faithfully, reversibly,  
and with full interpretability across discrete and continuous domains.

